\section{Introduction}

Guttation is among the most visible and least explained fungal behaviours. Droplets form on
mycelium across a wide taxonomic range, and their composition---where it has been
characterised---includes polyols, sugars, organic acids, proteins and secondary
metabolites. Proposed functions range from osmotic regulation and excretion of metabolic
by-products to deployment of extracellular enzymes and storage of defensive compounds. These
hypotheses make different predictions about the metabolism of the cells that produce the
droplets, but testing them requires knowing which cells those are.

Growing \textit{Pleurotus ostreatus} in a mycoponic ceramic-tube system
\cite{porterfield2026mycoponics} makes discrete mycelial structures accessible for individual
sampling, including a distinct exudate-producing structure we term the exudophore.
\textbf{[AUTHORS: one paragraph describing the exudophore morphologically and stating how it
is distinguished from exuding mycelium. This is required for the manuscript to stand.]}

Genome-scale metabolic models offer a way to move from a list of differentially expressed
genes to a statement about pathway operation. Constraint-based reconstructions integrate the
full stoichiometry of an organism's metabolism, and transcriptome-weighted contextualisation
methods extract from them a subnetwork consistent with a particular expression state. Applied
to tissues rather than conditions, these methods promise a metabolic account of
differentiation. We report both what that approach yielded and a caution about how its output
should be interpreted.

\section{Results}

\subsection{Context-specific models differ in size but not in the way that matters}

We reconstructed a genome-scale model for \textit{P.\ ostreatus} BOM\_ss5 (5,247 reactions,
2,287 able to carry flux on the defined culture medium; described in the companion resource
paper) and contextualised it separately against each tissue transcriptome using RIPTiDe \cite{jenior2020riptide}, which
weights reactions by transcript abundance and requires no expression threshold. All 881 model
genes were represented in the expression matrix.

The four models differ in size, and the exudophore retains the largest network: 124 reactions
against 109 for exuding mycelium, 100 for fuzzy mycelium and 93 for nodule
(Fig.~\ref{fig:models}a,b). Taken at face value this suggests a metabolically expansive
tissue, and 39 reactions appear in the exudophore model and in none of the other three.

That interpretation does not survive checking. Of those 39 reactions, \textbf{28 carry no gene
association whatsoever}---they are transport and exchange reactions whose presence reflects
what the network needs in order to carry flux, not what the tissue transcribes. Of the 11 that
do have gene support, most are expressed as highly or more highly in nodule than in the
exudophore: succinate:NAD\textsuperscript{+} oxidoreductase, pyruvate decarboxylase, fructose
1,6-bisphosphatase and ornithine ketoacid aminotransferase all fall into this category. Only
alcohol/methanol oxidase and the mannitol/sorbitol dehydrogenase pair are genuinely
exudophore-elevated.

Comparing the distribution of expression ratios for reactions unique to the exudophore model
against those shared with other models shows the two are barely distinguishable
(Fig.~\ref{fig:models}c). \textbf{Membership of a context-specific model is therefore not, by
itself, evidence of tissue-specific metabolism.} We report this because the reaction lists
such methods produce are commonly interpreted directly, and because the failure is not visible
without going back to the expression the model was built from.

\subsection{Direct expression testing across the network identifies 176 elevated reactions}

We therefore tested the evidence directly. Each model reaction was mapped through its
gene--protein--reaction rule to the genes that could catalyse it, and reaction-level expression
was summarised per tissue as the maximum across isoenzymes. Of 1,899 reactions with both gene
evidence and measurable expression, \textbf{176 are at least twofold higher in the exudophore
than in the highest of the other three tissues} (Fig.~\ref{fig:enzymes}a).

These converge on four functional themes.

\textbf{Polyol metabolism.} D-mannitol:NADP\textsuperscript{+} 2-oxidoreductase and D-sorbitol
2-dehydrogenase are elevated in the exudophore relative to every other tissue, but at
1.75-fold they fall below the twofold cut applied above. They are included here because they
are also among the small number of gene-supported reactions that the contextualised models
place in the exudophore and in no other tissue, and because of what mannitol is known to do. Mannitol is the
dominant soluble carbohydrate of many basidiomycetes and the principal osmolyte reported in
fungal exudate droplets. Its synthesis lowers water potential, which is the most parsimonious
physical mechanism for drawing water into a forming droplet.

\textbf{Oxidative one-carbon metabolism.} Alcohol/methanol oxidase shows the largest
expression contrast of any enzyme in the curated set, and aldehyde dehydrogenase is elevated
from a high baseline. Alcohol oxidase oxidises short-chain alcohols with molecular oxygen and
releases hydrogen peroxide; aldehyde dehydrogenase disposes of the resulting aldehyde. The
same gene supporting the alcohol oxidase reaction was independently identified as a
replicate-supported exudophore marker in the transcriptional analysis, by a method that makes
no use of the metabolic network.

\textbf{Cyanide detoxification.} Thiosulfate:cyanide sulfurtransferase (rhodanese) is elevated
approximately fourfold. The transcriptional analysis independently recovered cyanide hydratase
and cyanate hydratase among the replicate-supported markers. Two enzymatically distinct routes
for handling cyanide, recovered by two independent analyses, is more than coincidence deserves.

\textbf{Oxylipin synthesis.} Linoleate 8R-dioxygenase is among the most strongly elevated
reactions. Fungal oxylipins act as developmental and quorum signals, and their synthesis in a
structure that produces an extracellular liquid phase is suggestive.

\subsection{Flux comparison is uninformative on a rich medium}

We sampled flux distributions for each contextualised model and compared them. Differences
were small and dominated by amino-acid uptake: the mycoponic medium supplies peptone, malt
extract and tryptic soy broth, so biosynthetic pathways are bypassed by transport and the
models converge on similar solutions. Predicted biomass flux differed by roughly 2\% between
exudophore and baseline.

This is a limitation of the experimental system rather than of the method, and it is worth
stating for others attempting the same comparison: \textbf{constraint-based flux contrast
between tissues requires a medium restrictive enough to make biosynthesis load-bearing.} On a
rich undefined medium the informative signal lies in which reactions are expressed, not in how
much flux they carry.

\section{Discussion}

\subsection{A biochemical model of the exudophore}

The enzymes elevated in the exudophore compose a coherent scheme. Mannitol synthesis provides
the osmotic gradient that draws water into a droplet. Alcohol oxidase, acting on short-chain
alcohols, generates hydrogen peroxide in the extracellular space; aldehyde dehydrogenase and
formate dehydrogenase dispose of the aldehyde and formate produced, and cytochrome-c peroxidase
and peroxisomal catalase---both recovered independently---consume the peroxide. This is a
controlled oxidative cycle rather than uncontrolled leakage, and hydrogen peroxide is the
obligatory co-substrate of the class II peroxidases that mediate white-rot lignin attack.

The attraction of this scheme is that it makes the droplet functional rather than incidental: a
liquid phase in which oxidative chemistry can be conducted at a distance from the hypha, with
the osmolyte that creates the phase and the oxidase that charges it upregulated together.

\subsection{Predictions}

The scheme is testable without further sequencing.
\begin{enumerate}
\item Exudate droplets should be mannitol-rich, and mannitol concentration should exceed that
of the surrounding mycelium.
\item Droplets should contain hydrogen peroxide, and alcohol oxidase activity should be
detectable in droplet fluid.
\item Peroxidase activity in droplets should increase on addition of a suitable alcohol
substrate.
\item If oxylipin synthesis is genuinely elevated, 8R-hydroxy fatty acids should be detectable
in droplet fluid.
\end{enumerate}
The first two are straightforward assays and would together decide whether the scheme is right.

\subsection{Limitations}

The exudophore rests on two libraries after quality filtering, and all expression contrasts
carry that limitation. The metabolic model is a draft: its biomass objective is coarse and
uncurated, mass balance has not been curated, and three reactions were gapfilled on topological
rather than genomic evidence. Only 881 of 12,705 genes are represented in the network at all,
so any metabolic function without an EC-mapped homologue is invisible here---and lineage-specific
secreted proteins, the class most likely to matter in an exudate-producing structure, are
exactly the class least likely to be mapped. All enzyme assignments are inferred from homology;
none has been tested in this organism. The biochemical scheme above is a hypothesis assembled
from converging expression evidence, not a demonstrated pathway.

\section{Methods}

Culture, sequencing, read processing and the genome-scale reconstruction are described in the
companion resource paper; all code is in the accompanying repository.

\textbf{Contextualisation.} The gapfilled BOM\_ss5 model was contextualised separately against
each tissue with RIPTiDe \cite{jenior2020riptide}, using DESeq2 \cite{love2014deseq2} size-factor normalised counts for all 12,705 genes with
no expression filter, averaged within tissue. Using the \texttt{filterByExpr}-filtered matrix
instead pruned the models to roughly 90 reactions, because a gene absent from that matrix is
indistinguishable from a gene expressed at zero; the unfiltered matrix is the correct input.
500 flux samples were drawn per tissue. Model handling used COBRApy \cite{ebrahim2013cobrapy}; the reconstruction is built on the ModelSEED biochemistry database \cite{seaver2021modelseed}.

\textbf{Reaction-level expression.} Model reactions were mapped through their GPR rules to
locus tags via the \texttt{protein\_id} attribute of the stock GTF CDS records. Reaction
expression per tissue was taken as the maximum across isoenzymes, on the grounds that any one
suffices to catalyse the reaction. Reactions were ranked by the ratio of exudophore expression
to the highest of the other three tissues; a floor of 5 normalised counts was applied to avoid
ratios driven by near-zero denominators.

\noindent\textbf{[AUTHORS: sampling procedure and exudophore morphology, as above.]}

\clearpage
\section*{Figures}

\begin{figure}[htbp]\centering
\includegraphics[width=\textwidth]{fig6_tissue_models.pdf}
\caption{\textbf{Context-specific models, and why their reaction lists mislead.}
\textbf{a}, Reactions retained in each tissue model. \textbf{b}, How many tissue models each
reaction appears in. \textbf{c}, Expression ratio (exudophore over the highest other tissue)
for reactions unique to the exudophore model versus reactions shared with other models. The
distributions overlap: model membership does not imply elevated expression.}\label{fig:models}
\end{figure}

\begin{figure}[htbp]\centering
\includegraphics[width=\textwidth]{fig7_reaction_expression.pdf}
\caption{\textbf{Reaction-level expression across the metabolic network.}
\textbf{a}, Exudophore against the highest other tissue for all 1,899 gene-supported
reactions; dashed line marks twofold, orange points the 176 reactions exceeding it.
\textbf{b}, Expression of the enzymes composing the proposed biochemistry. Formate
dehydrogenase and cytochrome-c peroxidase are absent from this panel because they carry no
gene association in the reconstruction; they enter the argument from the contextualised models
and from the transcriptional marker analysis respectively.}\label{fig:enzymes}
\end{figure}

\clearpage
\section*{Tables}
\begin{table}[htbp]
\caption{Context-specific metabolic models extracted from the 5,247-reaction reconstruction by transcriptome-weighted contextualisation.}\label{tab:models}
\footnotesize
\begin{tabular}{@{}lrrrr@{}}
\toprule
Tissue & Libraries & Reactions & Metabolites & Genes \\
\midrule
Exuding mycelium & 3 & 109 & 122 & 77 \\
Fuzzy mycelium & 4 & 100 & 108 & 53 \\
Exudophore & 2 & 124 & 147 & 76 \\
Nodule & 3 & 93 & 109 & 53 \\
\bottomrule
\end{tabular}
\begin{tablenotes}[flushleft]\footnotesize\item All four were built from the same base model and the same 881 mapped genes.\end{tablenotes}
\end{table}

\begin{table}[htbp]
\caption{Enzymes composing the proposed exudophore biochemistry. Values are DESeq2-normalised expression averaged within tissue; ratio is exudophore over the highest other tissue.}\label{tab:enzymes}
\footnotesize
\begin{tabular}{@{}llrrrrrr@{}}
\toprule
Enzyme & EC & Reaction & Exudo. & Exuding & Fuzzy & Nodule & Ratio \\
\midrule
Alcohol/methanol oxidase & 1.1.3.13 & rxn00432 & 90.26 & 18.07 & 25.84 & 6.88 & 3.49 \\
Aldehyde dehydrogenase & 1.2.1.3 & rxn00506 & 175.58 & 23.5 & 15.92 & 46.93 & 3.74 \\
Mannitol 2-dehydrogenase & 1.1.1.138 & rxn20173 & 27.0 & 7.81 & 15.45 & 13.07 & 1.75 \\
Sorbitol 2-dehydrogenase & 1.1.1.14 & rxn24346 & 27.0 & 7.81 & 15.45 & 13.07 & 1.75 \\
Rhodanese & 2.8.1.1 & rxn01416 & 21.26 & 4.27 & 5.54 & 4.13 & 3.84 \\
Linoleate 8R-dioxygenase & 1.13.11.- & rxn04867 & 2.21 & 0.16 & 0.32 & 0.07 & 6.9 \\
\bottomrule
\end{tabular}
\begin{tablenotes}[flushleft]\footnotesize\item Formate dehydrogenase and cytochrome-c peroxidase are omitted: they carry no gene association in the reconstruction and so have no reaction-level expression.\end{tablenotes}
\end{table}

